\section*{GPT Outline for Introduction}

\begin{enumerate}
\item \textbf{Problem and Question.}

Start with calibration and calibeating under the quadratic/Brier score.
Ask the paper's central question directly: what survives if evaluation is
by an arbitrary proper Lipschitz scoring rule instead of Brier?

\item \textbf{Why This Is Nontrivial.}

Explain that Brier is special because it supports the familiar
decomposition into calibration plus refinement. Say that once we move to
general proper scoring rules, that decomposition is not automatic,
especially for fractional or continuous binning schemes.

\item \textbf{What the Paper Introduces.}

Introduce the notions of uniformly proper-calibrated and uniformly
proper-calibeating procedures over normalized $1$-Lipschitz proper
scoring rules. Briefly justify the Lipschitz restriction as the way to
make uniform comparison meaningful.

\item \textbf{Main Message: An Asymmetry.}

State the two headline facts early: calibration does lift to proper
calibration, whereas calibeating does not in general lift to proper
calibeating. This should be the conceptual center of the introduction.

\item \textbf{Background Literature.}

Use three short paragraphs:

\begin{itemize}
\item one on proper scoring rules: Savage, Schervish, Gneiting-Raftery,
and related entropy/divergence interpretations;
\item one on adversarial calibration: Dawid, Foster-Vohra, Foster,
smooth/continuous calibration in Foster-Hart, and related
approachability work;
\item one on calibeating and neighboring work: Foster-Hart,
Lee-Noarov-Pai-Roth, Chen-Huang-Jordan-Luo, and possibly Gupta-Ramdas
as an applied extension.
\end{itemize}

\item \textbf{Main Contributions.}

State the contributions in theorem order:

\begin{itemize}
\item exact decomposition for proper scoring rules when the binning
refines the forecast partition;
\item approximate decomposition for $\delta$-local fractional binnings;
\item every calibrated procedure is proper-calibrated;
\item an explicit counterexample showing that calibeating need not imply
proper calibeating;
\item positive proper-calibeating results for the simple online
averaging rule, for joint-binning stochastic procedures, and for
deterministic continuously proper-calibrated procedures.
\end{itemize}

\item \textbf{Why Joint Binning Matters.}

Give a short proof-level preview: the obstruction is not merely changing
the score, but losing the refinement structure when bins mix different
forecasts. Explain that the joint $\mathbf{b\times c}$ binning is what
restores the right decomposition logic.

\item \textbf{Relation to Broader Calibration Work.}

Add one short paragraph noting that calibration is currently being
extended in many directions, including multicalibration,
omniprediction, tail calibration, swap/KL calibration, and
arbitrary-sequence recalibration. Keep this brief; it is context, not
the center of the paper.

\item \textbf{Roadmap.}

End with a standard section-by-section map.
\end{enumerate}

\bigskip

\noindent \textbf{Editorial Recommendation.} If the introduction is kept
short, I would organize it around one sentence:

\begin{quote}
Calibration is universally robust across proper Lipschitz scoring rules,
but calibeating is not; this paper identifies exactly what additional
structure restores that robustness.
\end{quote}
